
Machine learning research increasingly relies on a narrow set of benchmark datasets, raising concerns that concentration may create epistemic lock-in---a self-reinforcing cycle constraining methodological diversity. We analyze 12,600 papers from three major ML venues (NeurIPS, ICML, ICLR) over seven years (2018--2024) using Herfindahl-Hirschman concentration indices derived from Papers With Code metadata. Citation networks strongly propagate benchmark standards: papers citing each other exhibit substantially higher dataset overlap than random pairs (Jaccard similarity 0.318 vs. 0.014, Cohen's d = 1.93). Prior-year concentration predicts individual paper adoption of standard benchmarks ($\beta$ = 56.75, p < 0.001). However, lagged panel regression reveals no temporal lock-in: the coefficient on prior-year concentration is positive rather than negative ($\beta$ = +1.60, p = 0.098), and Granger causality tests show zero significant venues. These results suggest that while the ML community transmits evaluation standards through scholarly networks, it co-evolves with its benchmarks rather than becoming deterministically trapped by them.

